\documentclass[10pt,a4paper]{amsart}
\usepackage{amsmath,amssymb,amsthm,graphicx}
% Retained minimal preamble for the Stacks Project chapter more-morphisms.tex (More on Morphisms).
% Verbatim excerpt of the chapter preamble: Xy-pic setup (native preamble lines 14-18) and the native
% theorem declarations with the mathematical macros (native preamble lines 154-217). Copied mechanically;
% nothing added, removed or redefined. The wrapper class and the amsmath/amssymb/amsthm/graphicx base
% packages are supplied by the assembler, so the class/externaldocument part of the native preamble is omitted.
% For commutative diagrams we use Xy-pic
\usepackage[all]{xy}

% We use 2cell for 2-commutative diagrams.
\xyoption{2cell}
% Theorem environments.
%
\theoremstyle{plain}
\newtheorem{theorem}[subsection]{Theorem}
\newtheorem{proposition}[subsection]{Proposition}
\newtheorem{lemma}[subsection]{Lemma}

\theoremstyle{definition}
\newtheorem{definition}[subsection]{Definition}
\newtheorem{example}[subsection]{Example}
\newtheorem{exercise}[subsection]{Exercise}
\newtheorem{situation}[subsection]{Situation}

\theoremstyle{remark}
\newtheorem{remark}[subsection]{Remark}
\newtheorem{remarks}[subsection]{Remarks}

\numberwithin{equation}{subsection}

% Macros
%
\def\lim{\mathop{\mathrm{lim}}\nolimits}
\def\colim{\mathop{\mathrm{colim}}\nolimits}
\def\Spec{\mathop{\mathrm{Spec}}}
\def\Hom{\mathop{\mathrm{Hom}}\nolimits}
\def\Ext{\mathop{\mathrm{Ext}}\nolimits}
\def\SheafHom{\mathop{\mathcal{H}\!\mathit{om}}\nolimits}
\def\SheafExt{\mathop{\mathcal{E}\!\mathit{xt}}\nolimits}
\def\Sch{\mathit{Sch}}
\def\Mor{\mathop{\mathrm{Mor}}\nolimits}
\def\Ob{\mathop{\mathrm{Ob}}\nolimits}
\def\Sh{\mathop{\mathit{Sh}}\nolimits}
\def\NL{\mathop{N\!L}\nolimits}
\def\CH{\mathop{\mathrm{CH}}\nolimits}
\def\proetale{{pro\text{-}\acute{e}tale}}
\def\etale{{\acute{e}tale}}
\def\QCoh{\mathit{QCoh}}
\def\Ker{\mathop{\mathrm{Ker}}}
\def\Im{\mathop{\mathrm{Im}}}
\def\Coker{\mathop{\mathrm{Coker}}}
\def\Coim{\mathop{\mathrm{Coim}}}

% Boxtimes
%
\DeclareMathSymbol{\boxtimes}{\mathbin}{AMSa}{"02}

%
% Macros for moduli stacks/spaces
%
\def\QCohstack{\mathcal{QC}\!\mathit{oh}}
\def\Cohstack{\mathcal{C}\!\mathit{oh}}
\def\Spacesstack{\mathcal{S}\!\mathit{paces}}
\def\Quotfunctor{\mathrm{Quot}}
\def\Hilbfunctor{\mathrm{Hilb}}
\def\Curvesstack{\mathcal{C}\!\mathit{urves}}
\def\Polarizedstack{\mathcal{P}\!\mathit{olarized}}
\def\Complexesstack{\mathcal{C}\!\mathit{omplexes}}
% \Pic is the operator that assigns to X its picard group, usage \Pic(X)
% \Picardstack_{X/B} denotes the Picard stack of X over B
% \Picardfunctor_{X/B} denotes the Picard functor of X over B
\def\Pic{\mathop{\mathrm{Pic}}\nolimits}
\def\Picardstack{\mathcal{P}\!\mathit{ic}}
\def\Picardfunctor{\mathrm{Pic}}
\def\Deformationcategory{\mathcal{D}\!\mathit{ef}}
\newenvironment{reference}{}{}
\newenvironment{slogan}{}{}
\newenvironment{history}{}{}
\usepackage{cleveref}
\crefname{originallemma}{original lemma}{original lemmas}
\Crefname{originallemma}{Original lemma}{Original lemmas}
\makeatletter
\newlabel{algebra-lemma-mod-injective-general}{{046Y}{source}{}{sourcepointer-1}{}}
\newlabel{algebra-lemma-mod-injective-general@cref}{{[originallemma][0][]046Y}{[1][0][]source}{}{}{}}
\makeatother
\crefname{originallemma}{original lemma}{original lemmas}
\Crefname{originallemma}{Original lemma}{Original lemmas}
\makeatletter
\newlabel{more-algebra-lemma-valuation-ring-flat-essentially-finite-type}{{0GSE}{source}{}{sourcepointer-2}{}}
\newlabel{more-algebra-lemma-valuation-ring-flat-essentially-finite-type@cref}{{[originallemma][0][]0GSE}{[1][0][]source}{}{}{}}
\makeatother
\title{Injectivity of a module map over a valuation ring is detected on the special fibre}
\author{The Stacks Project Authors}
\date{stacks-project:more-morphisms.tex:local-snapshot-20261011; GNU FDL 1.2 or later}
\begin{document}
\maketitle
\noindent\textit{Editorial note.} Statement span is lines 5043-5051 of sources/stacks-more-morphisms/pinned-original-source.tex; proof span is lines 5053-5091, the complete author proof (37 body lines).\par\medskip
\noindent\textit{Editorial note.} Navigation stubs: 2 cited-but-not-included labels are supplied as non-load-bearing pointers with their pinned tags where the snapshot tag index has them.\par\medskip
\noindent\textit{Editorial note.} extra\_packages is empty on purpose: the retained preamble already loads xy with options and a second declaration clashes in the assembler wrapper.\par\medskip
\noindent\textit{Editorial note.} Text is retained verbatim, including the author's own typos and placeholders.\par\medskip
\noindent\textit{Allowed standard prerequisites (editorial).}\begin{itemize}
\item Explicit prerequisite boundary. Valuation rings and their local homomorphisms, flatness criteria over a local ring, and essentially finitely type hypotheses; the reduction is retained.
\end{itemize}\medskip
\ifdefined\hypertarget\hypertarget{sourcepointer-1}{}\fi
\noindent\textit{Source reference.} Original lemma 046Y: sources/stacks-more-morphisms/pinned-original-source.tex, source lines 5091--5091. Cited by the retained statement or proof but not part of the collected spans; supplied only as a navigation pointer. The argument does not depend on its content.\par\medskip
\ifdefined\hypertarget\hypertarget{sourcepointer-2}{}\fi
\noindent\textit{Source reference.} Original lemma 0GSE: sources/stacks-more-morphisms/pinned-original-source.tex, source lines 5091--5091. Cited by the retained statement or proof but not part of the collected spans; supplied only as a navigation pointer. The argument does not depend on its content.\par\medskip
\begin{lemma}
\label{lemma-mod-injective-valuation-ring}
Let $A$ be a valuation ring. Let $A \to B$ is a local homomorphism of
local rings which is essentially of finite type.
Let $u : N \to M$ be a map of finite $B$-modules.
Assume $M$ is flat over $A$ and
$\overline{u} : N/\mathfrak m_A N \to M/\mathfrak m_A M$ is injective.
Then $u$ is injective and $M/u(N)$ is flat over $A$.
\end{lemma}

\begin{proof}
We will deduce this lemma from
Algebra, Lemma \ref{algebra-lemma-mod-injective-general}
(please note that we exchanged the roles of $M$ and $N$).
To do the reduction we will use More on Algebra, Lemma
\ref{more-algebra-lemma-valuation-ring-flat-essentially-finite-type}
to reduce to the finitely presented case.

\medskip\noindent
By assumption we can write $B$ as a quotient of the localization of
a polynomial algebra $P = A[x_1, \ldots, x_n]$ at a prime ideal
$\mathfrak q$. Then we can think of $u : N \to M$ as a map
of finite $P_\mathfrak q$-modules. Hence we may and do assume that
$B$ is essentially of finite presentation over $A$.

\medskip\noindent
Next, the $B$-module $N$ is finite but perhaps not of finite presentation.
Write $N = \colim N_\lambda$ as a filtered colimit of finitely presented
$B$-modules with surjective transition maps. For example choose
a presentation $0 \to K \to B^{\oplus r} \to N \to 0$,
write $K$ as the union of its finite submodules $K_\lambda$, and
set $N_\lambda = \Coker(K_\lambda \to B^{\oplus r})$.
The module $N/\mathfrak m_A N$ is of finite presentation
over the Noetherian ring $B/\mathfrak m_A B$.
Hence for $\lambda$ large enough we have
$N_\lambda/\mathfrak m_A N_\lambda = N/\mathfrak m_A N$.
Now, if we can show the lemma for the composition
$u_\lambda : N_\lambda \to M$, then we conclude that
$N_\lambda = N$ and the result holds for $u$.
Hence we may and do assume $N$ is of finite presentation over $B$.

\medskip\noindent
By More on Algebra, Lemma
\ref{more-algebra-lemma-valuation-ring-flat-essentially-finite-type}
the module $M$ is of finite presentation over $B$.
Thus all the assumptions of
Algebra, Lemma \ref{algebra-lemma-mod-injective-general}
hold and we conclude.
\end{proof}

\end{document}
